\section*{Your turn}

Check out tag \code{ch14} of the companion repository and run \code{npm ci}
once. Exercises 1 to 4 and 8 to 10 need the database, the broker, the seven
services and the router. Exercise 5 needs all of that too, plus a second
terminal and Reviews restarted on a different port. Exercises 6 and 7 need node
and nothing started at all.

There is a broker in the compose file now, so the line from
chapter~\ref{ch:hard-modeling-problems} grows one service:

\begin{minted}[fontsize=\footnotesize,breaklines]{text}
docker compose up -d --build mosaic-db mosaic-nats mosaic-catalog \
    mosaic-pricing mosaic-inventory mosaic-accounts mosaic-reviews \
    mosaic-ordering mosaic-nodes mosaic-router
docker compose ps
\end{minted}

Ten containers, all healthy. If \code{mosaic-nats} is the only one that will not
start, something else has 4222 or 8222.

\begin{enumerate}
  \item \textbf{Watch both subscriptions answer one write.} Run the harness
  the gate runs, with a product identifier and a customer who has not reviewed
  that product:

  \begin{minted}[fontsize=\footnotesize,breaklines]{text}
node scripts/subscription-run.mjs --product "..." --customer "..."
  \end{minted}

  Seven checks, all green at \code{ch14}. Read the two subprotocol lines first:
  the same offer, made to two servers, answered differently. Then stop
  \code{mosaic-nats} and run it again, and note which checks fail and which do
  not.

  \item \textbf{Make the graph pay for a field.} Open a subscription to
  \code{onReviewAdded} selecting \code{id} and \code{rating} only, submit three
  reviews, and count \code{Mosaic.RequestTimeline} lines in the Accounts
  console. Then do it again with \code{author \{ displayName \}} in the
  selection set. The difference is the running cost of one field, and the
  client chose it.

  \item \textbf{Read the plan for both subscriptions.} Send each of the two
  subscription fields to port 3002 with the query-plan headers
  chapter~\ref{ch:enter-the-router} used, and compare the trigger nodes. One
  names a subgraph and one does not. Then count the entity fetches under each
  and work out, before reading
  section~\ref{sec:ch14-the-subgraph-with-no-service} again, why the
  event-driven one needs an extra hop to fetch a field its own publisher
  already had in memory.

  \item \textbf{Publish an event nobody wrote.} With any NATS client, publish
  a message naming a review that exists:

  \begin{minted}[fontsize=\footnotesize,breaklines]{text}
subject  mosaic.reviewAdded.<the product's global identifier>
body     {"__typename":"Review","id":"<the review's global identifier>"}
  \end{minted}

  A client subscribed to \code{reviewPublished} on that product gets a full
  review that no mutation produced. Then send the same body with the type name
  removed. Note what the client is told, which service it appears to be about,
  and how long it takes you to connect that message to the edit you made.

  \item \textbf{Watch the handshake, then break it.} Start Reviews on another
  port and put the recorder in front of its usual one:

  \begin{minted}[fontsize=\footnotesize,breaklines]{text}
ASPNETCORE_URLS=http://localhost:5115 dotnet run --project src/Mosaic.Reviews -c Release
node scripts/subscription-wire.mjs 5105 5115
  \end{minted}

  Open a subscription through the router and read the \code{101}. Now delete
  the pinned subprotocol line from the reviews entry in
  \code{federation/mosaic.yaml}, recompose with \code{npm run compose}, and do
  it again. The subprotocol changes, and so do the names of two of the messages
  in the frames that follow.

  \item \textbf{Find out what the composer will let you say.} Print these in
  turn:

  \begin{minted}[fontsize=\footnotesize,breaklines]{text}
node scripts/realtime-cases.mjs --print documented-spelling-does-nothing
node scripts/realtime-cases.mjs --print unknown-subprotocol-is-ignored
node scripts/realtime-cases.mjs --print unknown-protocol-drops-the-key
  \end{minted}

  Three inputs a reviewer would pass and a composer accepts. Two of them end in
  a default and one ends in a missing key, which is a different thing. Decide
  which of the three you would rather find in production, and what would have
  had to run for you to find any of them.

  \item \textbf{Read a subgraph that is a file.} Open
  \code{schema/streams.graphql} beside the entry for it in
  \code{federation/mosaic.yaml}, then compose and find both halves of that
  entry in the execution config: the routing URL, which is in the subgraph list
  beside the other seven, and the eighth datasource, which has no url of any
  kind. Then satisfy yourself that nothing is listening on the port the routing
  URL names, and work out what that says about which of the two the router
  reads.

  \item \textbf{Break the transport on purpose.} Change the reviews
  \code{protocol} to \code{sse}, recompose, restart the router, and open a
  subscription. Now go looking for the cause: first in the subgraph console,
  then in the router log as it ships, then with \code{log\_level} in
  \code{router/config.yaml} turned up from \code{info} to \code{debug}. There
  is one line in the whole trail that names the problem, and it names it in a
  field rather than a message. Write down how long you would have spent on this
  without the recorder from exercise 5, then set the protocol to
  \code{sse\_post} and watch it work.

  \item \textbf{Ask for something in two parts.} Send the router a query with
  \code{@defer} on a fragment whose fields come from another subgraph, with
  \code{Accept: multipart/mixed}, and look at the response's content type. Then
  introspect the router for its directive list and introspect any subgraph for
  the same thing. Decide what you would tell a front-end team that had built a
  loading state around the directive the schema offers them.

  \item \textbf{Verify it in Postman.} The collection for this chapter is
  \code{postman/mosaic-realtime}, seven requests: six against the router and
  one against Reviews.

  \begin{minted}[fontsize=\footnotesize,breaklines]{text}
npx newman run postman/mosaic-realtime.postman_collection.json \
    -e postman/mosaic-realtime.local.postman_environment.json
  \end{minted}

  Twenty-three assertions, all green at \code{ch14}. The one request that does
  not go to the router is the one that shows the two subscription fields are
  not the same kind of thing. Note also what the collection does not test, and
  why decision 43 in this book's spec says a node script had to exist for that
  half.
\end{enumerate}

Exercise 4 is the one to do properly. Everything else in this chapter is a
property of a transport or a directive; that one is a property of trusting a
message. A subject on a broker is an interface with no schema behind it, and
the only thing standing between a typo and a graph that answers confidently
with the wrong thing is a convention two files agree about by hand.
Chapter~\ref{ch:automated-tests} is where a test suite gets told about
contracts nothing can compile.
